\appendix
\cleardoublepage
\phantomsection

\chapter{Apéndice A}\label{ap:A}

\section*{Detalle de relaciones críticas del modelo de datos}
Para complementar la descripción del modelo de datos, se detallan las relaciones de mayor impacto funcional:
\begin{itemize}
    \item \textbf{Usuario--Pedido (1:N):} un usuario puede registrar múltiples pedidos; cada pedido pertenece a un único cliente.
    \item \textbf{Pedido--OrderItem (1:N):} un pedido contiene uno o varios ítems con cantidad y precio unitario histórico.
    \item \textbf{Producto--OrderItem (1:N):} un producto puede aparecer en múltiples pedidos y conserva trazabilidad de ventas.
    \item \textbf{Usuario--Rol (N:M):} un usuario puede tener más de un rol (cliente, repartidor, administrador).
    \item \textbf{Producto--StockLocal--Local:} el inventario se gestiona por local para evitar inconsistencias de disponibilidad.
\end{itemize}

Estas relaciones sustentan la integridad transaccional del sistema, la trazabilidad de estados en el flujo de pedidos y la consistencia de inventario en tiempo real.


\chapter{Apéndice B}\label{ap:B}

\section*{Diagrama de red del sistema}
El diagrama de red ilustra la arquitectura general del sistema y la interacción entre usuarios web y móviles, la API y servicios externos.

\begin{figure}[htbp]
    \centering
    \resizebox{0.78\textwidth}{!}{% Redes: diagrama estilizado similar al resto de tikz del proyecto
\begin{tikzpicture}[node distance=10mm and 18mm, every node/.style={draw, rectangle, align=center, font=\small, rounded corners=2pt}]
  \node (usuarios) {Usuarios Web\\Cliente - Admin};
  \node (internet) [below=of usuarios] {Internet};
  \node (web) [below left=of internet, xshift=-6mm] {Frontend Web\\Navegador};
  \node (app) [below right=of internet, xshift=6mm] {App Móvil\\Android - iOS};
  \node (api) [below=24mm of internet] {API Backend\\REST};
  \node (db) [below left=of api] {Base de Datos};
  \node (chat) [below=of api] {Chatbot IA\\RAG};
  \node (ext) [below right=of api] {Servicios Externos\\Mapas - Pagos - WhatsApp};

  \draw[->, >=latex] (usuarios) -- (internet);
  \draw[->, >=latex] (internet) -- (web);
  \draw[->, >=latex] (internet) -- (app);
  \draw[->, >=latex] (web) -- (api);
  \draw[->, >=latex] (app) -- (api);
  \draw[->, >=latex] (api) -- (db);
  \draw[->, >=latex] (api) -- (chat);
  \draw[->, >=latex] (api) -- (ext);
\end{tikzpicture}
}
    \caption{Diagrama de red y arquitectura de la plataforma SaaS web y móvil}
    \label{fig:diagrama_red}
\end{figure}

\textit{Nota.} El diagrama de red muestra la interacción entre el frontend web y móvil, la API backend, la base de datos, el chatbot con inteligencia artificial integrado en la plataforma web y servicios externos como geolocalización; los medios de pago son efectivo y transferencia (sin pasarela de pagos).

\chapter{Apéndice C}\label{ap:C}

\section*{Procedimiento detallado de validación en \textit{staging}}
Este apéndice amplía el procedimiento de validación técnica descrito en el capítulo Método, con el detalle de los pasos ejecutados en entorno controlado.

\subsection*{1. Despliegue en entorno de \textit{staging}}
El backend se desplegó en un servidor VPS de Hetzner, utilizando Docker para contenerizar la aplicación y la base de datos PostgreSQL. Se configuraron ajustes de seguridad básicos (helmet, CORS, rate limiting) y variables de entorno mediante ConfigModule. La API se ejecutó en HTTPS y se activaron logs estructurados para monitoreo. El panel web se desplegó en Vercel y la aplicación móvil se distribuyó como APK para dispositivos de prueba.

\subsection*{2. Orientación a usuarios de prueba}
Se realizó una sesión de orientación breve (aproximadamente 1 hora) con participantes del escenario piloto simulado (administrador, personal de cocina y repartidores), explicando flujos principales y tareas de prueba.

\subsection*{3. Diseño de escenarios de prueba}
Se definieron escenarios predefinidos para cubrir flujos críticos: (1) registro de pedido completo por cliente, (2) confirmación y preparación en cocina, (3) asignación y entrega por repartidor, (4) gestión de productos y categorías por administrador, (5) consulta al chatbot en la plataforma web. Cada escenario utilizó datos simulados.

\subsection*{4. Ejecución de tareas guiadas}
Los usuarios ejecutaron tareas bajo observación estructurada, registrando tiempos, errores y completitud de flujos. Las pruebas se realizaron en sesiones controladas durante un período acotado.

\subsection*{5. Recolección y análisis de datos}
Se aplicó el cuestionario SUS al finalizar tareas guiadas y se consolidaron fichas de observación con métricas de desempeño. Se recopilaron logs para analizar disponibilidad, latencia y tasa de errores.

\begin{table}[htbp]
\centering
\caption{Tareas evaluadas en pruebas de usabilidad}
\label{tab:apendice-tareas-usabilidad}
\small
\begin{tabular}{|>{\raggedright\arraybackslash}p{0.8cm}|>{\raggedright\arraybackslash}p{4.0cm}|>{\raggedright\arraybackslash}p{2.4cm}|>{\raggedright\arraybackslash}p{5.2cm}|}
\hline
N.º & Tarea & Rol & Criterio de éxito \\
\hline
1 & Crear un nuevo pedido & Cliente & Completar en < 3 min \\
2 & Consultar estado de pedido & Cliente & Localizar en < 30 seg \\
3 & Agregar producto al catálogo & Administrador & Completar sin errores \\
4 & Asignar pedido a repartidor & Administrador & Completar en < 1 min \\
5 & Marcar pedido como entregado & Repartidor & Completar sin ayuda \\
6 & Realizar consulta al chatbot en la plataforma web & Cliente & Obtener respuesta pertinente \\
\hline
\end{tabular}
\end{table}

\subsection*{6. Respaldo de pruebas unitarias e integración}
Se documentó evidencia técnica de ejecución de pruebas en backend con Jest y validaciones de integración de endpoints críticos. El resumen de ejecución se presenta en la Tabla \ref{tab:apendice-evidencia-pruebas}; el detalle completo (salidas de consola, capturas y checklist de casos) se conservó en la bitácora técnica del proyecto.

\begin{table}[htbp]
\centering
\caption{Evidencia resumida de pruebas unitarias y de integración}
\label{tab:apendice-evidencia-pruebas}
\small
\begin{tabular}{|>{\raggedright\arraybackslash}p{2.0cm}|>{\raggedright\arraybackslash}p{4.0cm}|>{\raggedright\arraybackslash}p{2.5cm}|>{\raggedright\arraybackslash}p{4.2cm}|}
\hline
\textbf{Tipo} & \textbf{Módulo / alcance} & \textbf{Resultado} & \textbf{Evidencia registrada} \\
\hline
Unitaria & AuthService, OrdersService, ProductsService & Aprobadas & Salida Jest + trazabilidad por identificador de caso \\
\hline
Integración & Flujo login + creación de pedido + cambio de estado & Aprobadas & Registro de requests/responses y verificación de estados en BD \\
\hline
Integración & Gestión de inventario por pedido y validación de stock & Aprobadas & Comparación antes/después en tablas de stock y logs de servicio \\
\hline
Integración & Acceso por roles (admin, cliente, repartidor) & Aprobadas & Pruebas de autorización con respuestas 200/401/403 según rol \\
\hline
\end{tabular}
\end{table}

\chapter{Apéndice D}\label{ap:D}

\section*{Síntesis de la guía de entrevista semiestructurada (diagnóstico Fase 1)}

La guía de entrevista utilizada en el levantamiento inicial del dominio se estructura por bloques temáticos: (1) datos generales del negocio y antigüedad; (2) gestión de pedidos (canales actuales, volumen, tiempos, errores); (3) comunicación con cocina y repartidores; (4) control de inventarios; (5) atención al cliente y herramientas digitales en uso. Las preguntas son abiertas para profundizar en procesos operativos, necesidades y expectativas. Este apéndice presenta una síntesis del instrumento empleado como apoyo para la definición del escenario de referencia del estudio.

\chapter{Apéndice E}\label{ap:E}

\section*{Arquitectura detallada de la plataforma}
Este apéndice integra, en una sola vista textual, la arquitectura de la plataforma y las responsabilidades por capa para facilitar su revisión técnica sin depender únicamente de diagramas.

\begin{table}[htbp]
\centering
\caption{Detalle de arquitectura por capas y componentes}
\label{tab:apendice-arquitectura-detalle}
\small
\begin{tabular}{|>{\raggedright\arraybackslash}p{2.5cm}|>{\raggedright\arraybackslash}p{2.3cm}|>{\raggedright\arraybackslash}p{2.5cm}|>{\raggedright\arraybackslash}p{4.3cm}|}
\hline
\textbf{Capa} & \textbf{Componente} & \textbf{Tecnología} & \textbf{Responsabilidad principal} \\
\hline
Presentación & App móvil & React Native + Expo & Interfaz para cliente, cocina y repartidor; ejecución de flujos operativos. \\
\hline
Presentación & Panel web & Next.js + Tailwind CSS & Administración de catálogo, pedidos, repartidores y visualización de KPIs. \\
\hline
Negocio / API & Backend modular & NestJS + TypeScript & Reglas de negocio, seguridad JWT/RBAC, exposición de endpoints REST. \\
\hline
Datos & Persistencia transaccional & PostgreSQL (Supabase) + TypeORM & Gestión de entidades, relaciones e integridad de pedidos, usuarios e inventario. \\
\hline
Servicios externos & Mensajería y asistencia IA & Servicio web de chatbot RAG & Atención automatizada y consultas frecuentes contextualizadas. \\
\hline
Infraestructura & Despliegue & VPS (Hetzner), Vercel, Docker & Operación en \textit{staging}, publicación y continuidad técnica del prototipo. \\
\hline
\end{tabular}
\end{table}

La arquitectura se diseñó para mantener separación de responsabilidades, bajo acoplamiento entre capas y escalabilidad incremental, permitiendo evolucionar módulos sin afectar de forma directa el resto de componentes.

